\scanpage{154}
\mspage{78}
% archive sheet number: 153
% L1–3
avec la compatibilité entre (67) (68) :
\[
\begin{aligned}
&\text{si }\vec a\in\vec A\text{ est incidente (\uline{avec orientations induites})}\\
&\text{à }\vec f\in\vec F,\quad\text{alors }\vec\tau_A(\vec a)=\vec\tau_F(\vec f)
\end{aligned}\tag*{(69)}
\]
% L4–8
en d'autres termes on a commutativité dans
\diagram{\[
\begin{tikzcd}[ampersand replacement=\&]
\&\mathrm{Rep}(S,F)\arrow[dl,"\mathrm{can}"']\arrow[dr,"\mathrm{can}"]\&\\
\vec A\arrow[dr,"\vec\tau_A"']\&\&\vec F\arrow[dl,"\vec\tau_F"]\\
\&\unsure{\vec\Pi|\vec\pi}\&\ .
\end{tikzcd}\tag*{(70)}
\]}

% L9–13
Je voudrais maintenant dévisser un peu plus la situation du diagramme (60). Fixons nous un $s_0\in S$, et considérons son stabilisateur \deleted{dans} $G_{s_0}$ dans $G$, on a donc, \margin{\unsure{dévissant|divisant}}
\[
1\mathrel{\unsure{-|\longrightarrow}}\underset{\substack{\wr\Vert\ \text{non canoniquement}\\\mathbb Z/n\mathbb Z}}{G^+_{0s_0}}
 \longrightarrow G_{s_0}\longrightarrow(\mathbb Z/n\mathbb Z)^*
\]
% L14–17
je dis que $G_{s_0}\mathrel{\overset{\chi\mid G_{s_0}}{\unsure{-|\longrightarrow}}}(\mathbb Z/n\mathbb Z)^*$ est épi -- c'est formel, à partir du fait que $G_0^+$ est transitif sur l'ensemble $S$, donc on a une structure d'extension canonique :
% L18–19
\[
1\longrightarrow\underset{\substack{\wr\Vert\\\mathbb Z/n\mathbb Z\ \text{\uline{non} can.}}}{G^+_{0s_0}}
\longrightarrow G_{s_0}\longrightarrow(\mathbb Z/n\mathbb Z)^*\longrightarrow1\ .
\tag*{(71)}
\]
% slanted
\margin{On a envie de noter ce groupe $U_{s_0}$ (``\unsure{unipotent|unipotents}'' !)}
% slanted
\margin{on a envie de le noter $B'_{\unsure{s_0}}$ (``Borel'') / mais il lui manque encore un facteur $G_m$ pour le mériter\ldots}
